% TOP_PROBLEM: {"id": 8700024, "title": "Conjecture 3.6 — Roy", "category_id": 3, "category": {"id": 3, "name": "graph_theory", "display_name": "Graph Theory", "description": "Problems involving graphs, networks, and their properties.", "slug": "graph-theory", "order_index": 3, "created_at": "2026-07-31T15:26:25.671Z"}, "status": "open", "problem_number": "AMR-086-0024", "set_id": 13}
% FIRST_POSTED: 2026-10-01
% ENTRY_KIND: statement_only
% UnsolvedMath ID 8700024; source catalogue code AMR-086-0024
% !TeX program = lualatex
% Definition and sources only; this file is not an LLM solution attempt.
\documentclass[11pt,a4paper]{article}
\usepackage[margin=25mm]{geometry}
\usepackage{fontspec}
\setmainfont{lmroman10-regular.otf}[BoldFont=lmroman10-bold.otf,
 ItalicFont=lmroman10-italic.otf,BoldItalicFont=lmroman10-bolditalic.otf]
\usepackage{amsmath,amssymb,mathtools}
\usepackage{unicode-math}
\setmathfont{latinmodern-math.otf}
\AtBeginDocument{\let\setminus\smallsetminus}
\usepackage{xurl,hyperref}
\hypersetup{colorlinks=true,urlcolor=blue}
\setcounter{secnumdepth}{0}
\setlength{\parindent}{0pt}
\setlength{\parskip}{5pt}
\setlength{\emergencystretch}{3em}
\begin{document}
\section*{Conjecture 3.6 — Roy}\label{8700024}
{\small UnsolvedMath ID 8700024 · AMR-086-0024 · Graph Theory · AMR Open Problem Lists}
\subsection{Definitions and mathematical statement}
Define the set of complex logarithms of algebraic numbers by
\[
\mathcal L=\{z\in\mathbb C:e^z\in\overline{\mathbb Q}^{\times}\}.
\]
All logarithm branches are included, and zero belongs to this set. Let $M=(m_{ij})$ be a $4\times4$ matrix with entries in $\mathcal L$, satisfying $M^T=-M$. In particular its diagonal entries are zero. Matrix rank and column span below are taken over $\mathbb C$.

Suppose $\operatorname{rank}_{\mathbb C}M\le2$. Roy's conjecture asks whether at least one of the following alternatives must hold:
\[
\exists\,0\ne q\in\mathbb Q^4:\ q^TM=0,
\qquad\text{or}\qquad
\operatorname{im}_{\mathbb C}M\cap\mathbb Q^4\ne\{0\}.
\]
The first alternative means a nontrivial rational linear dependence among the four row vectors, regarded as vectors in $\mathbb C^4$. The second means that there exist complex coefficients combining the columns into a nonzero vector whose four coordinates are all rational.

The two fields play different roles: rank at most two already implies complex row dependence, whereas the first conclusion demands dependence with rational coefficients. Likewise, the second conclusion concerns the full complex column space, not just rational combinations of columns.

The assertion includes zero matrices and all other permitted ranks, with no further nonvanishing assumption on the entries. The problem is to prove this disjunction for every matrix satisfying the logarithm and skew-symmetry hypotheses, or produce a matrix for which both rational conclusions fail.
\subsection{Short English statement}
Must a low-rank skew-symmetric matrix of logarithms of algebraic numbers have a rational row relation or a nonzero rational vector in its column space?
\subsection{Sources}
\begin{itemize}
\item[S1] ULAM.AI and the UnsolvedMath Contributors, supplied problems.json, record 8700024 (AMR-086-0024), AMR Open Problem Lists. Catalogue identity and source transcription; CC BY 4.0.
\url{https://huggingface.co/datasets/ulamai/UnsolvedMath}.
\item[S2] Waldschmidt - Open Diophantine Problems (2004), Conjecture 3.6. Original source indexed by AMR-086-0024.
\url{https://arxiv.org/abs/math/0312440}.
\item[S3] M. Waldschmidt, Open Diophantine Problems, Moscow Mathematical Journal 4 (2004), 245–305. Author PDF supplies the surrounding definitions and hypotheses.
\url{https://webusers.imj-prg.fr/~michel.waldschmidt/articles/pdf/odp.pdf}.
\end{itemize}
\end{document}
